\documentclass[12pt,a4paper]{report}
% Compile with XeLaTeX in TeXstudio.
\usepackage{fontspec}
\usepackage{polyglossia}
\setmainlanguage{mongolian}
\setmainfont{Times New Roman}

\usepackage[a4paper,top=2cm,bottom=2cm,left=2.5cm,right=2cm]{geometry}
\usepackage{setspace}
\usepackage{graphicx}
\usepackage{float}
\usepackage{amsmath}
\usepackage{array}
\usepackage{tabularx}
\usepackage{booktabs}
\usepackage{xurl}
\usepackage[hidelinks]{hyperref}
\usepackage{listings}
\usepackage{xcolor}
\usepackage{longtable}
\newcolumntype{P}[1]{>{\raggedright\arraybackslash}p{#1}}

\onehalfspacing
\setlength{\parindent}{1.12cm}
\setlength{\parskip}{2pt}
\renewcommand{\contentsname}{ГАРЧИГ}
\renewcommand{\figurename}{Зураг}
\renewcommand{\tablename}{Хүснэгт}
\renewcommand{\bibname}{Ашигласан материал}

\lstset{
	basicstyle=\ttfamily\small,
	breaklines=true,
	frame=single,
	columns=fullflexible,
	keepspaces=true
}

% ---- Макронууд ----
\newcommand{\repo}{\url{https://github.com/Naidan2006/sudoku-validator}}
\newcommand{\shot}[3]{\begin{figure}[H]\centering\includegraphics[width=0.85\textwidth]{#1}\caption{#2}\label{#3}\end{figure}}
% 9x9 хүснэгт: 9 мөр, мөр бүр "a&b&c&d&e&f&g&h&i" хэлбэртэй
\newcommand{\sgrid}[9]{{\small\setlength{\tabcolsep}{3.5pt}\begin{tabular}{|ccc|ccc|ccc|}\hline #1\\#2\\#3\\\hline #4\\#5\\#6\\\hline #7\\#8\\#9\\\hline\end{tabular}}}

\lstset{basicstyle=\ttfamily\small,breaklines=true,frame=single}

\begin{document}
	
	% ---------------- НҮҮР ХУУДАС ----------------
	\begin{titlepage}
		\begin{center}
			{\large ШИНЖЛЭХ УХААН, ТЕХНОЛОГИЙН ИХ СУРГУУЛЬ\par}
			\vspace{0.2cm}
			{\normalsize Сургууль: Мэдээлэл холбоо технологийн их сургууль\par}
			\vspace{2cm}
				\IfFileExists{shutis_logo.png}{%
				\includegraphics[height=3.2cm,keepaspectratio]{shutis_logo.png}\par
				\vspace{0.35cm}
			}{}
			{\Huge\bfseries БИЕ ДААЛТ 1\par}
			\vspace{0.8cm}
			{\Large\bfseries F.CSM302 ҮЙЛДЛИЙН СИСТЕМ\par}
			\vspace{0.4cm}
			{\large\bfseries SUDOKU SOLUTION VALIDATOR\par}
			\vspace{0.3cm}
			{\normalsize Олон урсгалт Sudoku шийдэл шалгах програм\par}
		\end{center}
		
		\vspace{1cm}
		\begin{flushleft}
		{\large\bfseries Хуваарилалтын код:} {\normalsize II-4-5\par}
		{\large\bfseries Сурах бичгийн төсөл:} {\normalsize Operating System Concepts (2018) Chapter 4: Threads & Concurrency, Programming Project 5\par}
		{\large\bfseries Гүйцэтгэсэн оюутан / оюутны код:\par}\begin{flushleft} 
											{\normalsize О.Энхжаргал B232230049\par}\\
										   	{\normalsize А.Найдан B241930004\par}\\
											{\normalsize Д.Алтан-Очир B221940001\par}\\
											{\normalsize Б.Аззаяа B242270803\par}\\
											{\normalsize Н.Мөнхбаяр B242270058\par}\\
											{\normalsize Ц.Тэргэл B241940128\par}\end{flushleft}
		{\large\bfseries Шалгасан багш:} {\normalsize Г.Ганбат\par}
		{\large\bfseries GitHub repository:} \url{https://github.com/Naidan2006/sudoku-validator}	
		\end{flushleft}
		
		\begin{center}
			Улаанбаатар хот\\
			2026 он
		\end{center}
	\end{titlepage}
	\pagenumbering{roman}
	
	% ---------------- ХУРААНГУЙ ----------------
	\chapter*{Хураангуй}
	\addcontentsline{toc}{chapter}{Хураангуй}
	
	Энэхүү төсөл нь Үйлдлийн системийн хичээлийн хүрээнд C програмчлалын хэл болон POSIX Threads (\texttt{pthread}) ашиглан 9$\times$9 хэмжээтэй Sudoku хүснэгтийн зөв эсэхийг шалгах програм боловсруулахад чиглэнэ. Програмын шалгалт нь мөр, багана болон есөн ширхэг 3$\times$3 дэд хүснэгтийг тус тус шалгах 11 урсгалын зохион байгуулалтад суурилна. Төслийн repository-д терминалаар ажиллах хувилбар (\texttt{main.c}) болон GTK3 график интерфэйстэй хувилбар (\texttt{gui.c}) багтсан. Шалгалтын үндсэн логик нь \texttt{sudoku\_core.c}, нийтлэг зарлалууд нь \texttt{sudoku\_core.h} файлд байрлана. Хэрэглэгчийн өгөгдөлд хоосон нүд, 1--9-ийн хүрээнээс гадуурх утга, мөр/багана/дэд хүснэгтийн давхардал байгаа эсэхийг шалгаж, хэсэгчилсэн үр дүнг \texttt{PASS}/\texttt{FAIL}, эцсийн төлөвийг \texttt{VALID}/\texttt{INVALID} хэлбэрээр харуулах зорилготой. Туршилтын аргачлалд зөв Sudoku, хоосон нүд, буруу утга, давхардсан тоо болон интерфэйсийн цэвэрлэх үйлдлийг хамруулна. 
	
	\noindent\textbf{Түлхүүр үг:} үйлдлийн систем, олон урсгал, pthread, Sudoku Validator, GTK3.
	
	\tableofcontents
	\clearpage
	\pagenumbering{arabic}
	
	% ---------------- 1 УДИРТГАЛ ----------------
	\chapter{Удиртгал}
	\section{Зорилго}
	C хэл болон \texttt{pthread} ашиглан Sudoku хүснэгтийн мөр, багана, 3$\times$3 дэд хүснэгтүүдийг олон урсгалаар шалгадаг, үр дүнг нэгтгэн харуулдаг програмыг судлах, хэрэгжүүлэх, турших.
	
	\section{Зорилтууд}
	\begin{itemize}
		\item Sudoku шалгах дүрэм болон олон урсгалын үндсэн ойлголтыг судлах.
		\item Шалгах ажлыг 11 урсгалд хуваарилах зохион байгуулалтыг тайлбарлах.
		\item Консолын болон GTK3 график интерфэйстэй хувилбарын бүтцийг тодорхойлох.
		\item Зөв болон буруу өгөгдөлд зориулсан туршилтын тохиолдол боловсруулах.
		\item Туршилтын бодит үр дүнг бүртгэж, шаардлагатай нотолгоог screenshot болон хүснэгтээр үзүүлэх.
	\end{itemize}
	
	\section{Ач холбогдол}
	Энэхүү төсөл нь үйлдлийн системийн хичээлээр үздэг урсгал үүсгэх, ажлыг хэсэгчлэн хуваарилах, урсгалуудын гүйцэтгэсэн шалгалтын үр дүнг нэгтгэх ойлголтыг жижиг, ойлгомжтой асуудал дээр харуулна. График интерфэйс нь шалгалтын төлөвүүдийг хэрэглэгчид харагдах байдлаар үзүүлэх боломж олгоно.
	
	\section{Хамрах хүрээ}
	Repository-ийн тайлбарт 9$\times$9 Sudoku шалгах, мөр/багана/дэд хүснэгтийг шалгах, \texttt{PASS}/\texttt{FAIL} болон \texttt{VALID}/\texttt{INVALID} үр дүн харуулах боломж дурдагдсан. Мөн терминалын хувилбар болон GTK3 интерфэйстэй хувилбар байна.
	
	% ---------------- 2 ШААРДЛАГА ----------------
	\chapter{Сурах бичгийн шаардлага}
	
	Энэхүү төсөл нь Silberschatz, Galvin, Gagne нарын Operating System Concepts(2018) сурах бичгийн 4-р бүлэг Threads & Concurrency-н Programming Project 5 буюу Sudoku Solution Validator юм. Хуваарилалтын код:II-4-5. Сурах бичгийн төсөл нь 9x9 Sudoku шийдлийг шалгахдаа бүх мөрийг шалгах нэг урсгал, бүх баганыг шалгах нэг урсгал, есөн 3x3 дэд хүснэгтийг тус бүр шалгах есөн урсгал, нийт 11 урсгал ашиглахыг шаарддаг.
	
	\section{Шаардлага ба хэрэгжилтийн зураглал}
	\begin{table}[H]
		\centering
		\caption{Шаардлага ба төслийн хэрэгжилт}
		\small
		\begin{tabularx}{\textwidth}{|p{0.32\textwidth}|X|}
			\hline
			\textbf{Шаардлага} & \textbf{Repository-ийн баримтад дурдсан хэрэгжилт}\\
			\hline
			9$\times$9 Sudoku хүснэгтийг шалгах & 81 нүдтэй хүснэгтийн дүрмийн шалгах; GUI хувилбарт хүснэгтийн нүдэнд цифр оруулах.\\
			\hline
			Мөр, багана, дэд хүснэгт шалгах & Нэг урсгал мөрүүд, нэг урсгал баганууд, есөн урсгал тус бүр нэг 3$\times$3 дэд хүснэгтийг шалгана.\\
			\hline
			Шалгалтын төлөв харуулах & Хэсэг тус бүрийн үр дүнг \texttt{PASS}/\texttt{FAIL}, эцсийн үр дүнг \texttt{VALID}/\texttt{INVALID} гэж харуулна.\\
			\hline
			Оролтын алдаа боловсруулах & README-д хоосон нүд болон 1--9-ийн хүрээнээс гадуурх утгыг илрүүлэхийг дурдсан. Бодит зан төлөвийг туршилтаар баталгаажуулна.\\
			\hline
			Дахин шалгах боломж & Терминалын хувилбар дахин Sudoku шалгах боломжтой; GUI-д хүснэгтийг цэвэрлэж дахин шалгах үйлдэл тайлбарлагдсан.\\
			\hline
		\end{tabularx}
	\end{table}
	
	% ---------------- 3 ОНОЛ ----------------
	\chapter{Онолын үндэс}
	
	\section{Процесс ба урсгал}
	Процесс нь ажиллаж буй програмын тусгаарлагдсан гүйцэтгэлийн орчин бол thread нь тухайн процесс доторх гүйцэтгэлийн урсгал юм. Нэг процессын thread-үүд код болон процессын олон нөөцийг хуваалцаж болох боловч тус бүр өөрийн гүйцэтгэлийн төлөв, стек зэрэгтэй. Иймээс олон урсгал нь нэг асуудлын өөр өөр хэсгийг зэрэг ажиллуулах боломж олгохын зэрэгцээ shared data-д хандах болон үр дүнг зөв нэгтгэх шаардлага үүсгэнэ.
	
	\section{POSIX Threads (\texttt{pthread})}
	\texttt{pthread} нь POSIX орчинд урсгал үүсгэх, удирдах зориулалттай интерфэйс юм. Энэ төсөлд \texttt{pthread} нь Sudoku-ийн өөр өөр хэсгийг шалгах ажлуудыг тусдаа урсгалд гүйцэтгэхэд ашиглагдана. Төслийн Windows орчинд MSYS2 UCRT64 болон GCC ашиглахаар README-д заасан.
	
	\section{Sudoku шалгах алгоритм}
	Мөр, багана болон дэд хүснэгт тус бүрийн шалгалт нь тухайн 9 нүдэнд 1--9 цифр давхардалгүй байгаа эсэхийг тогтооно. Хэрэв шалгалт нь бүх цифр 1--9 яг нэг удаа байгааг баталдаг бол тухайн хэсэг PASS болно. Хэрэв хоосон нүдийг тусдаа алдаа гэж үздэг бол хоосон нүдний шалгалтыг дүрмийн шалгалтаас өмнө эсвэл нийт дүгнэлтэд нэмэлт нөхцөл болгон хэрэгжүүлэх шаардлагатай.
	
	\section{Урсгалын хуваарилалт ба өгөгдөл хуваалцах}
	Програм 11 шалгах урсгал ашиглана: нэг нь мөрүүд, нэг нь баганууд, ес нь тус тусын 3$\times$3 дэд хүснэгтийг шалгана. Урсгал бүр хүснэгтийг зөвхөн унших бөгөөд үр дүнгээ \texttt{results[]} массивын өөрийн тусдаа элементэд бичдэг (мөр: \texttt{results[0]}, багана: \texttt{results[1]}, дэд хүснэгт: \texttt{results[index+2]}) тул хоёр урсгал нэг нүдэнд бичихгүй. Иймээс уралдаан (race condition) үүсэхгүй бөгөөд mutex шаардлагагүй. Үндсэн урсгал бүх урсгалыг \texttt{pthread\_join()}-оор хүлээсний дараа л үр дүнг нэгтгэнэ. Шалгалтын явцад хүснэгтийн өгөгдлийг өөрчлөхгүй байхаар GUI талыг зохион байгуулсан (5.5 хэсэг).
	
	% ---------------- 4 ЗОХИОМЖ ----------------
	\chapter{Системийн зохиомж}
	
	\section{Компонентууд}
	\begin{itemize}
		\item \texttt{main.c} -- терминал дээр ажиллах хувилбар.
		\item \texttt{gui.c} -- GTK3 график интерфэйс.
		\item \texttt{sudoku\_core.c} -- Sudoku шалгах алгоритм болон урсгалын төлөвтэй холбоотой логик.
		\item \texttt{sudoku\_core.h} -- нийтлэг тодорхойлолт, функцийн зарлал.
		\item \texttt{tests/} -- тестийн тохиолдлуудын тайлбар.
		\item \texttt{docs/} -- гарын авлага болон дэлгэцийн зураг.
		\item \texttt{.github/workflows/} -- GitHub Actions хөрвүүлэлтийн тохиргоо.
	\end{itemize}
	
	\section{Өгөгдлийн урсгал}
	\begin{figure}[H]
		\centering
		\fbox{\parbox{0.88\textwidth}{
				\centering
				Хэрэглэгчийн 9$\times$9 Sudoku өгөгдөл\\[0.2cm]
				$\Downarrow$\\[0.2cm]
				Консол (\texttt{main.c}) эсвэл GTK3 интерфэйс (\texttt{gui.c})\\[0.2cm]
				$\Downarrow$\\[0.2cm]
				Sudoku шалгах үндсэн логик (\texttt{sudoku\_core.c})\\[0.2cm]
				$\Downarrow$\\[0.2cm]
				11 урсгал: мөр + багана + 9 дэд хүснэгт\\[0.2cm]
				$\Downarrow$\\[0.2cm]
				Хэсэг тус бүрийн \texttt{PASS}/\texttt{FAIL}\\[0.2cm]
				$\Downarrow$\\[0.2cm]
				Эцсийн төлөв: \texttt{VALID}/\texttt{INVALID}
		}}
		\caption{Sudoku Validator системийн өгөгдлийн урсгалын загвар}
	\end{figure}
	
	% ---------------- 5 ХЭРЭГЖИЛТ ----------------
	\chapter{Хэрэгжилт}
	
	\section{Хөгжүүлэлтийн орчин}
	Repository-ийн README-д дараах технологийг дурдсан:
	\begin{itemize}
		\item Програмчлалын хэл: C.
		\item Олон урсгал: POSIX Threads (\texttt{pthread}).
		\item График интерфэйс: GTK3.
		\item Хөрвүүлэгч: GCC.
		\item Орчин: Windows болон MSYS2 UCRT64.
		\item GUI build-д: \texttt{pkg-config} болон GTK3-ийн хөгжүүлэлтийн сангууд.
	\end{itemize}
	
	\section{Консолын хувилбарыг хөрвүүлэх ба ажиллуулах}
	README-д заасан командууд:
	\begin{lstlisting}
		gcc -Wall -Wextra -pthread main.c -o sudoku_validator.exe
		./sudoku_validator.exe
	\end{lstlisting}
	Консолын хувилбарт \texttt{Row 1}--\texttt{Row 9} хүртэл мөр бүрд зайгаар тусгаарласан 9 бүхэл тоо оруулна. Ажиллуулсан орчин болон гарсан бодит үр дүнг туршилтын бүртгэлд тэмдэглэнэ.
	
	\section{GTK3 график интерфэйсийг хөрвүүлэх ба ажиллуулах}
	README-д заасан командууд:
	\begin{lstlisting}
		gcc gui.c sudoku_core.c -o sudoku_gui.exe -pthread \
		$(pkg-config --cflags --libs gtk+-3.0)
		./sudoku_gui.exe
	\end{lstlisting}
	Энэ хувилбар ажиллахын тулд GTK3 сан болон шаардлагатай хөгжүүлэлтийн хэрэгслүүд MSYS2 UCRT64 орчинд суусан байх ёстой. GTK3 dependency байхгүй орчинд GUI build шууд ажиллана гэж үзэхгүй.
	
	\section{График интерфэйсийн үндсэн үйлдлүүд}
	Repository-ийн README-д тайлбарласнаар хэрэглэгч 81 нүдэнд 1--9 хүртэлх цифр оруулж, \textbf{Шалгах} товчоор үр дүнг харна. \textbf{Цэвэрлэх} товч нь хүснэгт болон өмнөх үр дүнг арилгах, \textbf{Заавар} товч нь ашиглах зааврыг нээх зориулалттай. Товч тус бүрийн бодит ажиллагааг GUI дээр ажиллуулж баталгаажуулна.
	
	\section{Кодын гол хэсгүүд}
	Кодын бүхэл файлыг тайланд хуулж оруулахын оронд \texttt{sudoku\_core.c} доторх урсгал үүсгэх, мөр/багана/дэд хүснэгт шалгах, үр дүнг нэгтгэх функцуудын бодит хэсгийг сонгон оруулах нь тохиромжтой. Энэ тайланд эх кодын бүрэн агуулга хавсаргагдаагүй тул функцийн нэр болон мөрийн дугаарыг таамгаар заагаагүй. Эцсийн хувилбарт бодит кодтой тулган богино кодын жишээ нэмнэ.
	
	% ---------------- 6 UI ТУРШИЛТ ----------------
	\chapter{Туршилт ба UI дүрслэл}
	
	\section{Багшаас оноосон үндсэн дүрслэл}
	Хичээлийн ерөнхий зааварт (Бие даалт 1, нийтлэг шаардлагын 4-р заалт) багшаас төсөлд оноосон үндсэн дүрслэл бүхий UI-ийг хөгжүүлж, бодит ажиллагаатай холбохыг шаардсан. GTK3 интерфэйс нь 9$\times$9 Sudoku хүснэгт, хэрэглэгчийн үйлдлийн товчлуурууд, шалгалтын үр дүн болон 11 шалгах урсгалын төлвийг харуулах бөгөөд PASS/FAIL, VALID/INVALID төлөв нь хэрэглэгчийн оруулсан бодит өгөгдлийг шалгасан үр дүнтэй холбогдсон.
	\section{Туршилтын орчин}
	Үйлдлийн систем: Windows 11; хөгжүүлэлтийн орчин: MSYS2 UCRT64; хөрвүүлэгч: GCC; урсгалын сан: POSIX Threads; график интерфэйс: GTK3; эх код: \repo. Нарийн хувилбарыг 5.2 хэсгийн хүснэгтээс харна уу.
	\section{Юуг өөрчилж турших вэ}
	Хэрэглэгч дараахыг өөрчилж туршиж болно: (1) 81 нүдний аль ч утга; (2) нүдийг хоосон үлдээх; (3) 1--9-өөс өөр тэмдэгт оруулах; (4) үр дүнг үзсэний дараа утгыг засаж дахин шалгах; (5) цэвэрлэх товчоор хүснэгтийг дахин эхлүүлэх. Ажиглах зүйл: хэсэг бүрийн PASS/FAIL, 11 урсгалын төлөв, эцсийн VALID/INVALID, мөн алдааны мэдэгдэл.
	\section{Хоосон нүдтэй оролт}
	Зөв хүснэгтийн (1,1) нүдийг хоосон үлдээсэн (Хүснэгт~\ref{fig:g1}-ийн тест 1). Програм шалгалт эхлүүлэхгүйгээр хоосон нүдний мөр, баганыг мэдэгдэх ёстой.\begin{figure}[H]
		\centering
		\includegraphics[width=0.85\textwidth]
		{Зураг 6.1.png}
		\caption{Хоосон нүдтэй оролтын алдааны мэдээпэп}
		\label{fig:valid-sudoku}
	\end{figure}
	\section{Буруу Sudoku}
	Зөв хүснэгтийн (1,2) нүдийг 3-аас 5 болгож өөрчилсөн (тест 2). Мөр 1, багана 2, дэд хүснэгт 1 FAIL, эцсийн төлөв INVALID байх ёстой.\begin{figure}[H]
		\centering
		\includegraphics[width=0.85\textwidth]
		{Зураг 6.2.png}
		\caption{Буруу Sudoku-ийн шалгалтын үр дүн}
		\label{fig:valid-sudoku}
	\end{figure}
	\section{Хэсэгчилсэн алдаа}
	Хоёр тусдаа алдаатай хүснэгтэд (тест 3) алдаатай хэсгүүд FAIL, бусад нь PASS байж, мөр, багана, дэд хүснэгтийн төлөв тус тусдаа харагдах ёстой.\begin{figure}[H]
		\centering
		\includegraphics[width=0.85\textwidth]
		{Зураг 6.3.png}
		\caption{Мөр, багана болон дэд хүснэгтийн шалгалтын үр дүн}
		\label{fig:valid-sudoku}
	\end{figure}
	\section{Зөв Sudoku}
	Зөв хүснэгтэд (тест 4) мөр, багана, есөн дэд хүснэгт бүгд PASS, эцсийн төлөв VALID байх ёстой. \begin{figure}[H]
		\centering
		\includegraphics[width=0.85\textwidth]
		{Зураг 6.4.png}
		\caption{Зөв Sudoku-ийн шалгалтын үр дүн}
		\label{fig:valid-sudoku}
	\end{figure}
	\section{Урсгалын төлөвийн самбар ба responsive интерфэйс}
	Самбар нь 11 урсгал тус бүрийн төлвийг харуулна. Зөв хүснэгтэд бүгд PASS болсон байх ёстой. \begin{figure}[H]
		\centering
		\includegraphics[width=0.85\textwidth]
		{Зураг 6.5.png}
		\caption{11 урсгалын төлөв болон шалгалтын үр дүнн}
		\label{fig:valid-sudoku}
	\end{figure}
	%
	
	% ---------------- 7 ТУРШИЛТЫН АРГА ----------------
\chapter{Туршилтын арга}
\section{Арга ба хүлээгдэх үр дүн}
Тест бүрд оролт, хүлээгдэж буй үр дүнг урьдчилан тодорхойлж, програмын бодит үр дүнтэй харьцуулсан. Дэд хүснэгтийг мөр-мөрөөр зүүнээс баруун тийш 1--9 гэж дугаарлана (дэд хүснэгт 1 нь зүүн дээд, 9 нь баруун доод). Нүдийг (мөр, багана) гэж заана. Тестийн оролтын хүснэгтүүдийг 7.2 хэсэгт бүрэн өгсөн тул өөр хүн ижил оролтоор давтаж шалгах боломжтой.
\begin{longtable}{P{0.06\textwidth}P{0.2\textwidth}P{0.32\textwidth}P{0.3\textwidth}}\caption{Туршилтын кейсүүд}\label{tab:tests}\\\toprule
	\textbf{№}&\textbf{Тест}&\textbf{Оролт}&\textbf{Хүлээгдэж буй үр дүн}\\\midrule\endfirsthead
	\toprule\textbf{№}&\textbf{Тест}&\textbf{Оролт}&\textbf{Хүлээгдэж буй үр дүн}\\\midrule\endhead
	1&Хоосон нүд&Зөв хүснэгт, (1,1) хоосон&``1-р мөр, 1-р багана хоосон байна'' гэсэн алдааны мэдэгдэл\\
	2&Буруу Sudoku&Зөв хүснэгт, (1,2)=5&Мөр 1, багана 2, дэд хүснэгт 1 FAIL; бусад PASS; INVALID\\
	3&Хэсэгчилсэн алдаа&(1,1)=3 ба (9,9)=8&Мөр 1,9; багана 1,9; дэд хүснэгт 1,9 FAIL; бусад PASS; INVALID\\
	4&Зөв Sudoku&Зөв хүснэгт&11 урсгал бүгд PASS; VALID\\
	5&Зөвхөн мөр алдаатай&(1,1)ба(2,1)-ийг солисон&Мөр 1,2 FAIL; багана, дэд хүснэгт бүгд PASS; INVALID\\
	6&Зөвхөн багана алдаатай&(1,1)ба(1,2)-ийг солисон&Багана 1,2 FAIL; мөр, дэд хүснэгт бүгд PASS; INVALID\\
	7&Зөвхөн дэд хүснэгт алдаатай&Мөр 1 ба мөр 4-ийг солисон&Дэд хүснэгт 1--6 FAIL, 7--9 PASS; мөр, багана бүгд PASS; INVALID\\
	8&Буруу тэмдэгт&Нүдэнд 0, 10, -1, a оруулах&0: ``1-р мөр, 1-р баганад 1--9 оруулна''; a, -: ижил алдаа; 10: нүд 1 тэмдэгтээр хязгаарлагдаж зөвхөн 1 үлдэнэ\\
	9&Засаж дахин шалгах&Тест 2-ын дараа (1,2)-ийг 3 болгох&Дахин шалгахад бүх хэсэг PASS; VALID\\
	10&Цэвэрлэх (reset)&Аль ч тестийн дараа Clear товч&Бүх нүд хоосон болж, статус ``Шалгаагүй'' болно; шалгалт явж байхад товч блоклогдоно\\
	\bottomrule\end{longtable}
\section{Тестийн өгөгдөл}
Дараах хүснэгтүүдэд тест бүрийн оролтыг үзүүлэв. Эхний хүснэгт (Зөв) нь жишиг бөгөөд нүд хоосон байвал нүд хоосон гэж үзнэ; өөрчилсөн нүдийг тодруулав.
\begin{figure}[H]\centering
	\begin{minipage}[t]{0.32\textwidth}\centering Тест 4 (Зөв)\par\sgrid{5&3&4&6&7&8&9&1&2}{6&7&2&1&9&5&3&4&8}{1&9&8&3&4&2&5&6&7}{8&5&9&7&6&1&4&2&3}{4&2&6&8&5&3&7&9&1}{7&1&3&9&2&4&8&5&6}{9&6&1&5&3&7&2&8&4}{2&8&7&4&1&9&6&3&5}{3&4&5&2&8&6&1&7&9}\end{minipage}\hfill
	\begin{minipage}[t]{0.32\textwidth}\centering Тест 1 (Хоосон нүд)\par\sgrid{ &3&4&6&7&8&9&1&2}{6&7&2&1&9&5&3&4&8}{1&9&8&3&4&2&5&6&7}{8&5&9&7&6&1&4&2&3}{4&2&6&8&5&3&7&9&1}{7&1&3&9&2&4&8&5&6}{9&6&1&5&3&7&2&8&4}{2&8&7&4&1&9&6&3&5}{3&4&5&2&8&6&1&7&9}\end{minipage}\hfill
	\begin{minipage}[t]{0.32\textwidth}\centering Тест 2 (Буруу)\par\sgrid{5&\textbf{5}&4&6&7&8&9&1&2}{6&7&2&1&9&5&3&4&8}{1&9&8&3&4&2&5&6&7}{8&5&9&7&6&1&4&2&3}{4&2&6&8&5&3&7&9&1}{7&1&3&9&2&4&8&5&6}{9&6&1&5&3&7&2&8&4}{2&8&7&4&1&9&6&3&5}{3&4&5&2&8&6&1&7&9}\end{minipage}
	\caption{Тест 4, 1, 2-ийн оролт}\label{fig:g1}
\end{figure}
\begin{figure}[H]\centering
	\begin{minipage}[t]{0.32\textwidth}\centering Тест 3 (Хэсэгчилсэн)\par\sgrid{\textbf{3}&3&4&6&7&8&9&1&2}{6&7&2&1&9&5&3&4&8}{1&9&8&3&4&2&5&6&7}{8&5&9&7&6&1&4&2&3}{4&2&6&8&5&3&7&9&1}{7&1&3&9&2&4&8&5&6}{9&6&1&5&3&7&2&8&4}{2&8&7&4&1&9&6&3&5}{3&4&5&2&8&6&1&7&\textbf{8}}\end{minipage}\hfill
	\begin{minipage}[t]{0.32\textwidth}\centering Тест 5 (зөвхөн мөр)\par\sgrid{\textbf{6}&3&4&6&7&8&9&1&2}{\textbf{5}&7&2&1&9&5&3&4&8}{1&9&8&3&4&2&5&6&7}{8&5&9&7&6&1&4&2&3}{4&2&6&8&5&3&7&9&1}{7&1&3&9&2&4&8&5&6}{9&6&1&5&3&7&2&8&4}{2&8&7&4&1&9&6&3&5}{3&4&5&2&8&6&1&7&9}\end{minipage}\hfill
	\begin{minipage}[t]{0.32\textwidth}\centering Тест 6 (зөвхөн багана)\par\sgrid{\textbf{3}&\textbf{5}&4&6&7&8&9&1&2}{6&7&2&1&9&5&3&4&8}{1&9&8&3&4&2&5&6&7}{8&5&9&7&6&1&4&2&3}{4&2&6&8&5&3&7&9&1}{7&1&3&9&2&4&8&5&6}{9&6&1&5&3&7&2&8&4}{2&8&7&4&1&9&6&3&5}{3&4&5&2&8&6&1&7&9}\end{minipage}
	\caption{Тест 3, 5, 6-ын оролт}\label{fig:g2}
\end{figure}
\begin{figure}[H]\centering
	\begin{minipage}[t]{0.4\textwidth}\centering Тест 7 (зөвхөн дэд хүснэгт)\par\sgrid{\textbf{8}&\textbf{5}&\textbf{9}&\textbf{7}&\textbf{6}&\textbf{1}&\textbf{4}&\textbf{2}&\textbf{3}}{6&7&2&1&9&5&3&4&8}{1&9&8&3&4&2&5&6&7}{\textbf{5}&\textbf{3}&\textbf{4}&\textbf{6}&\textbf{7}&\textbf{8}&\textbf{9}&\textbf{1}&\textbf{2}}{4&2&6&8&5&3&7&9&1}{7&1&3&9&2&4&8&5&6}{9&6&1&5&3&7&2&8&4}{2&8&7&4&1&9&6&3&5}{3&4&5&2&8&6&1&7&9}\end{minipage}
	\caption{Тест 7-ийн оролт (мөр 1 ба 4 солигдсон)}\label{fig:g3}
\end{figure}
Тест 8--10-д тусгай хүснэгт хэрэггүй: тест 8-д дурын нүдэнд 0, 10, $-1$, \texttt{a} утга оруулж, тест 9-д тест 2-ын хүснэгтийн (1,2) нүдийг 3 болгож, тест 10-д аль ч хүснэгтийн дараа цэвэрлэх товч дарсан.

%
	% ---------------- 8 ҮР ДҮН ----------------
	\chapter{Үр дүн}
	
	\section{Үр дүнг нэгтгэх}
	Туршилт бүрийн дараа бодит үр дүнг хүснэгтэд бүртгэж, хүлээгдсэн үр дүнтэй харьцуулна. Зөв Sudoku-д бүх шалгалт \texttt{PASS}, эцсийн төлөв \texttt{VALID} байх нь хүлээгдэнэ. Дүрэм зөрчсөн Sudoku-д холбогдох шалгалт \texttt{FAIL}, эцсийн төлөв \texttt{INVALID} байх нь хүлээгдэнэ. Хоосон нүд, буруу утга болон цэвэрлэх үйлдлийн бодит хариу үйлдлийг програм дээр ажиллуулж баталгаажуулах шаардлагатай.
	
	\section{Нотолгоо}
	Үр дүнгийн нотолгоонд дараах материалыг хавсаргана:
	\begin{itemize}
		\item Дөрвөн GUI screenshot.
		\item Тест бүрийн яг ашигласан 9$\times$9 өгөгдөл.
		\item Консолын гаралт эсвэл GUI дээрх бодит үр дүн.
		\item Туршилтын орчны хувилбарын мэдээлэл.
	\end{itemize}
	Энэ тайланг боловсруулах үед бодит туршилтын лог болон screenshot файлуудыг энд хавсаргаагүй тул үр дүнгийн тоон дүгнэлт, бодит PASS/FAIL тоог зориудаар зохиож оруулаагүй.
	\section{Гүйцэтгэлийн хэмжилт}
	Урсгалуудын ажиллах хугацаа хэмжигдээгүй тул гүйцэтгэлийн хурд, хурдатгал, CPU хэрэглээний тоон үр дүн тайланд оруулаагүй. Цаашид хэмжилт хийх бол ижил өгөгдөл, ижил орчинд олон удаагийн давталт хийж, хэмжилтийн арга болон үр дүнг тусад нь тайлагнах шаардлагатай.
	
	% ---------------- 9 ХЭЛЭЛЦҮҮЛЭГ ----------------
	\chapter{Хэлэлцүүлэг}
	Олон урсгал ашигласнаар Sudoku-ийн бие даасан хэсгүүдийг тус тусад нь шалгах боломжтой болсон. Тест 5--7 нь алдааг мөр, багана, дэд хүснэгтийн аль нэгэнд тусгаарлан үүсгэдэг тул PASS/FAIL үр дүн алдаа илэрсэн хэсгийг яг заадаг эсэхийг шалгана. GTK3 интерфэйс нь хүснэгт оруулах, шалгах, засварлан дахин шалгах боломж олгож, урсгалын төлөвийн самбараар 11 шалгалтын үр дүнг тус тусад нь харуулдаг.
	
	\textbf{Хязгаарлалт.}
	\begin{itemize}
		\item Програм Sudoku-г бодож шийдэхгүй, зөвхөн өгөгдсөн шийдлийн зөв эсэхийг шалгана.
		\item График хувилбарт GTK3 болон MSYS2 UCRT64 орчин шаардлагатай; Linux, macOS дээр туршаагүй.
		\item Гүйцэтгэлийн хурдыг хэмжээгүй тул 11 урсгалын хурдны давуу талын тухай тоон дүгнэлт хийгээгүй. 9$\times$9 хүснэгт жижиг тул урсгал үүсгэх зардал нь шалгалтын зардлаас их байж болно.
	\end{itemize}
	% ---------------- 10 ДҮГНЭЛТ ----------------
\chapter{Дүгнэлт}
Энэ төслөөр C хэл, POSIX Threads болон GTK3 ашигласан олон урсгалт Sudoku шалгагчийн консол болон график хувилбарыг боловсруулав. Шалгалтыг мөр, багана болон есөн 3$\times$3 дэд хүснэгтэд хуваарилж, нийт 11 урсгал ашигласан нь сурах бичгийн шаардлагатай нийцнэ. Хоосон нүд, буруу, хэсэгчилсэн алдаатай болон зөв Sudoku-г туршиж, үр дүнг screenshot-оор нотолсон. Цаашид оролтын баталгаажуулалтыг сайжруулах, Linux орчны заавар нэмэх, гүйцэтгэлийг хэмжих боломжтой.

%==================== ЛАВЛАГАА ====================
\chapter*{Ашигласан материал}\addcontentsline{toc}{chapter}{Ашигласан материал}
\begin{enumerate}
	\item A. Silberschatz, P. B. Galvin, G. Gagne. \textit{Operating System Concepts}, 10th ed. Wiley, 2018. Chapter 4, Programming Project 5.
	\item F.CSM302 Үйлдлийн систем, Бие даалт 1-ын заавар. Хичээлийн материал.
	\item The Open Group. POSIX Threads API. \url{https://pubs.opengroup.org/onlinepubs/9699919799/}
	\item GTK Project. GTK 3 Reference Manual. \url{https://docs.gtk.org/gtk3/}
	\item GNU Compiler Collection. \url{https://gcc.gnu.org/}
	\item Sudoku Validator төслийн эх код. \repo
	\item Код, тайлбар, баримтжуулалт боловсруулахад ChatGPT, тайлангийн LaTeX засварт Claude ашигласан.
\end{enumerate}
%==================== ХАВСРАЛТ ====================
\appendix
\chapter{Ажлын төлөвлөгөө ба хуваарилалт}
\begin{table}[H]\centering\caption{Төслийн ерөнхий мэдээлэл}\begin{tabular}{P{0.28\textwidth}P{0.62\textwidth}}\toprule
		\textbf{Талбар}&\textbf{Мэдээлэл}\\\midrule
		Хуваарилалтын код&II-4-5\\
		Төсөл&Sudoku Solution Validator\\
		Repository&\repo\\
		Хэл/сан/орчин&C, pthread, GTK3, GCC, MSYS2 UCRT64\\
		Үндсэн дүрслэл&11 урсгалын төлвийн самбар, PASS/FAIL, VALID/INVALID\\\bottomrule\end{tabular}\end{table}
\begin{longtable}{P{0.3\textwidth}P{0.4\textwidth}P{0.2\textwidth}}
	\caption{Багийн ажлын хуваарилалт}\\
	\toprule
	\textbf{Гишүүн} & \textbf{Хариуцсан ажил} & \textbf{Commit/file}\\
	\midrule\endfirsthead
	\toprule\textbf{Гишүүн} & \textbf{Хариуцсан ажил} & \textbf{Commit/file}\\\midrule\endhead
	Д.~Алтан-Очир (B221940001) & Sudoku шалгах алгоритмын шаардлага, мөр болон багана шалгах логикийг хянах; алгоритмын тайлбар. & \texttt{sudoku\_core.c}\\
	Б.~Аззаяа (B242270803) & GTK3 интерфэйсийн зохиомж, хэрэглэгчийн оролт болон товчлууруудын ажиллагааг шалгах. & \texttt{gui.c}\\
	О.~Энхжаргал (B232230049) & Олон урсгалын зохицуулалт, \texttt{pthread} функцүүдийн тайлбар болон эх кодын хяналт. & \texttt{sudoku\_core.c/.h}\\
	Н.~Мөнхбаяр (B242270058) & Туршилтын кейс бэлтгэх, зөв/буруу Sudoku-ийн үр дүнг шалгах, тестийн хүснэгт бүрдүүлэх. & тестийн оролт (тайлан 7.2)\\
	Ц.~Тэргэл (B241940128) & Хоосон нүд болон алдаатай оролтын туршилт, дэлгэцийн зураг болон хэрэглэгчийн зааврын хяналт. & \texttt{docs/screenshots/}\\
	А.~Найдан (B241930004) & GitHub repository, README/Manual-ийн уялдаа, тайланг нэгтгэх, LaTeX хувилбар, эцсийн шалгалт. & \texttt{README.md}, \texttt{main.c}, тайлан\\
	\bottomrule
\end{longtable}
Commit-үүдийг хоёр нэрээр (\texttt{Naidan2006}, \texttt{Naidan}; ижил email) бүртгэсэн тул гишүүн тус бүрийн commit-ийн тоог тусад нь ялгах боломжгүй.	
\end{document}
